\documentclass[11pt]{article}
\usepackage[margin=1in]{geometry}
\usepackage{amsmath,amssymb,amsthm}
\usepackage{enumitem}
\usepackage[hidelinks]{hyperref}

\newtheorem{theorem}{Theorem}
\newtheorem{lemma}{Lemma}
\newtheorem{definition}{Definition}
\newtheorem{remark}{Remark}

\newcommand{\linf}{\ell_\infty}
\newcommand{\card}[1]{\left|#1\right|}
\newcommand{\D}{D_\kappa}
\newcommand{\CR}{\operatorname{CR}}
\newcommand{\dens}{\operatorname{dens}}

\begin{document}

\begin{center}
{\Large A negative answer to the cardinal condensation-rank problem for
\(\ell_\infty(\Gamma)\)}

\vspace{0.4em}
{\large Packet for arXiv:1104.4896}
\end{center}

\section*{Source problem}

Gazor's paper \emph{Condensation rank of injective Banach spaces}
(arXiv:1104.4896) defines the usual condensation derivative and proves a
lower-bound theorem for infinite-dimensional injective Banach spaces.  At the
end of Section 3 the author asks the following cardinal analogue: after
defining a \(\kappa\)-condensation derivative, if
\(\kappa=2^{\card{\Gamma}}\), is the \(\kappa\)-condensation rank of
\(\ell_\infty(\Gamma)\) the first ordinal of cardinality \(\kappa\)?

The natural nontrivial convention is that a point is a
\(\kappa\)-condensation point of \(A\) when every neighbourhood of the point
meets \(A\) in at least \(\kappa\) points.  Under this convention, the answer
is no: the rank is one cardinal step larger.

\section*{Definitions}

Let \(\kappa\) be an infinite cardinal and \(X\) a topological space.  For
\(A\subseteq X\), put
\[
D_\kappa(A)=\{x\in X:\card{U\cap A}\geq \kappa
\text{ for every neighbourhood }U\text{ of }x\}.
\]
The transfinite iterates are defined by
\[
D_\kappa^0(A)=A,\qquad
D_\kappa^{\alpha+1}(A)=D_\kappa(D_\kappa^\alpha(A)),\qquad
D_\kappa^\lambda(A)=\bigcap_{\alpha<\lambda}D_\kappa^\alpha(A)
\]
for limit ordinals \(\lambda\).  Let \(\rho_\kappa(A)\) be the least ordinal
\(\alpha\) such that \(D_\kappa^\alpha(A)=D_\kappa^{\alpha+1}(A)\), and set
\[
\CR_\kappa(X)=\sup\{\rho_\kappa(A):A\subseteq X\}.
\]
For a cardinal \(\mu\), let \(\omega_\mu\) denote the initial ordinal of
cardinality \(\mu\).  Thus the ordinal asked for in the source problem is
\(\omega_\kappa\), while \(\omega_{\kappa^+}\) is the first ordinal whose
cardinality is strictly larger than \(\kappa\).

\section*{Proof intuition}

The upper bound is purely cardinal: along a strictly decreasing derivative
chain one can choose a different point lost at each successor step, so a space
of cardinality \(\kappa\) cannot support \(\kappa^+\) many strict losses.

The lower bound is a transfinite rank-tree construction inside
\(\ell_\infty(\Gamma)\).  The cube \(\{0,1\}^{\Gamma}\) gives
\(\kappa=2^{\card{\Gamma}}\) many pairwise distance-one directions.  These
directions let us place \(\kappa\) many small copies of any previously built
rank block near a root, in countably many shrinking shells.  The root then
becomes a \(\kappa\)-condensation point of every earlier derivative, survives
one derivative longer than the copies, and disappears immediately after that.
Limit ordinals are handled by separated sums of cofinally many earlier rank
blocks.

\section*{Main theorem}

\begin{theorem}
Let \(\Gamma\) be an infinite set and let
\(\kappa=2^{\card{\Gamma}}\).  With the convention above,
\[
\CR_\kappa(\ell_\infty(\Gamma))=\omega_{\kappa^+}.
\]
Consequently the equality proposed in the final problem of arXiv:1104.4896,
\(\CR_\kappa(\ell_\infty(\Gamma))=\omega_\kappa\), is false.
\end{theorem}

\begin{proof}
First note that
\[
\card{\ell_\infty(\Gamma)}=\kappa.
\]
Indeed, \(\{0,1\}^{\Gamma}\subseteq \ell_\infty(\Gamma)\) has cardinal
\(2^{\card{\Gamma}}\), while the set of all scalar functions on \(\Gamma\) has
cardinal
\[
(2^{\aleph_0})^{\card{\Gamma}}=2^{\card{\Gamma}}
\]
because \(\Gamma\) is infinite.

\medskip
\noindent\emph{Upper bound.}
Let \(A\subseteq \ell_\infty(\Gamma)\).  If
\(D_\kappa^\alpha(A)\supsetneq D_\kappa^{\alpha+1}(A)\) for every
\(\alpha<\lambda\), then we may choose
\[
x_\alpha\in D_\kappa^\alpha(A)\setminus D_\kappa^{\alpha+1}(A)
\qquad(\alpha<\lambda).
\]
These points are distinct: if \(\alpha<\beta\), then
\[
x_\beta\in D_\kappa^\beta(A)\subseteq D_\kappa^{\alpha+1}(A),
\]
whereas \(x_\alpha\notin D_\kappa^{\alpha+1}(A)\).  Hence
\(\card{\lambda}\leq \card{A}\leq \kappa\).  Therefore no derivative chain in
\(\ell_\infty(\Gamma)\) is strictly decreasing through all ordinals below
\(\omega_{\kappa^+}\), and
\[
\CR_\kappa(\ell_\infty(\Gamma))\leq \omega_{\kappa^+}.
\]

\medskip
\noindent\emph{Two elementary placement facts.}
Let
\[
V=\{0,1\}^{\Gamma}\setminus\{0\}\subseteq \ell_\infty(\Gamma).
\]
Then \(\card{V}=\kappa\), \(\|v\|=1\) for every \(v\in V\), and distinct
points of \(V\) have distance \(1\).

We shall use two standard consequences of this separated family.

First, if \(\{B_i:i\in I\}\) is a family of closed subsets of
\(\ell_\infty(\Gamma)\), \(\card{I}\leq\kappa\), and each \(B_i\) has
cardinality at most \(\kappa\), then scaled translates of the \(B_i\)'s can be
placed in pairwise uniformly separated balls inside any prescribed nonempty
ball.  For such a uniformly separated union,
\[
D_\kappa^\alpha\Big(\bigcup_i B_i\Big)=\bigcup_i D_\kappa^\alpha(B_i)
\quad\text{for every ordinal }\alpha. \tag{1}
\]
The reason is local: every point has a neighbourhood meeting at most one
component.

Second, suppose \(\{B_i:i\in I\}\) is a family of closed sets, indexed by
\[
I=\mathbb N\times\kappa,
\]
and suppose each \(B_i\) has cardinality at most \(\kappa\).  Given a point
\(p\) and a radius \(r>0\), choose numbers \(a_n\downarrow 0\) so fast that
the shells remain separated, for instance \(a_n=r2^{-2n-4}\).  Put the
components in balls of radius \(a_n/100\) around points \(p+a_nv_\xi\), where
\(\xi<\kappa\) and \(v_\xi\in V\).  The resulting closed set
\[
F=\{p\}\cup \bigcup_{(n,\xi)\in I} B_{n,\xi}
\]
has only one possible cross-component \(\kappa\)-condensation point, namely
\(p\).  Away from \(p\), a small neighbourhood meets only finitely many shells
and, in each such shell, at most one component.  Every neighbourhood of \(p\),
however, contains all components in sufficiently small shells and therefore
contains \(\kappa\) many components.

\medskip
\noindent\emph{Rank blocks.}
We prove by transfinite induction that for every nonzero
\(\alpha<\omega_{\kappa^+}\) and every nonempty ball \(B\) in
\(\ell_\infty(\Gamma)\), there is a closed set \(S_\alpha\subseteq B\), with
\(\card{S_\alpha}\leq\kappa\), such that
\[
D_\kappa^\beta(S_\alpha)\neq\varnothing\quad(\beta<\alpha),
\qquad
D_\kappa^\alpha(S_\alpha)=\varnothing. \tag{2}
\]
Thus \(S_\alpha\) has \(\kappa\)-condensation height exactly \(\alpha\).

For \(\alpha=1\), take a singleton.

Assume first that \(\alpha\) is a nonzero limit ordinal.  Since
\(\alpha<\omega_{\kappa^+}\), \(\card{\alpha}\leq \kappa\).  Choose a cofinal
set \(C\subseteq \alpha\) with \(\card{C}\leq\kappa\).  By the induction
hypothesis, for each \(\gamma\in C\) choose a closed rank block of height
\(\gamma\), and place these blocks as a uniformly separated family inside
\(B\).  Equation (1) shows that the derivative of the union is the union of
the derivatives.  Given \(\beta<\alpha\), some \(\gamma\in C\) satisfies
\(\beta<\gamma\), so the \(\beta\)-th derivative of the union is nonempty.
At stage \(\alpha\), every component has already disappeared.  Hence (2)
holds.

Now suppose \(\alpha=\beta+1\).  If \(\beta=0\), this is again the singleton
case.  If \(\beta>0\), choose for every \((n,\xi)\in\mathbb N\times\kappa\)
a closed block \(B_{n,\xi}\) of height \(\beta\), scaled small enough to fit
in the fan construction above, and set
\[
S_{\beta+1}=\{p\}\cup\bigcup_{(n,\xi)\in\mathbb N\times\kappa}B_{n,\xi}.
\]
For every \(\eta<\beta\), each component has nonempty
\(D_\kappa^\eta\), and every neighbourhood of \(p\) contains \(\kappa\) many
components.  Hence \(p\in D_\kappa^{\eta+1}(S_{\beta+1})\), and consequently
\(p\in D_\kappa^\beta(S_{\beta+1})\).  The components have empty
\(\beta\)-th derivative, so the fan observation gives
\[
D_\kappa^\beta(S_{\beta+1})=\{p\}.
\]
The singleton \(\{p\}\) has empty \(\kappa\)-condensation derivative, so
\[
D_\kappa^{\beta+1}(S_{\beta+1})=\varnothing.
\]
Moreover the same argument gives nonemptiness at every earlier stage.  This
completes the induction.

Since rank blocks exist for every nonzero
\(\alpha<\omega_{\kappa^+}\), the supremum of the ranks of subsets of
\(\ell_\infty(\Gamma)\) is at least \(\omega_{\kappa^+}\).  Together with the
upper bound, this proves
\[
\CR_\kappa(\ell_\infty(\Gamma))=\omega_{\kappa^+}.
\]
\end{proof}

\section*{Remarks on conventions}

The source problem is meaningful only under an ``at least \(\kappa\)''
interpretation of \(\kappa\)-condensation.  If one instead required strictly
more than \(\kappa\) points in every neighbourhood while also taking
\(\kappa=2^{\card{\Gamma}}=\card{\ell_\infty(\Gamma)}\), the first derivative
would be empty for every subset, so the proposed equality would fail
trivially.  The theorem above addresses the nontrivial convention.

\section*{Novelty and verification}

The local index check found no previous run packet for arXiv:1104.4896.  A
bounded web search for the phrases ``aleph-condensation'', ``condensation
rank'', and ``\( \ell_\infty(\Gamma) \)'' found only the source arXiv paper,
not a later answer to the final problem.  The proof uses only elementary
cardinal arithmetic, transfinite induction, and the explicit separated
family \(\{0,1\}^{\Gamma}\subseteq\ell_\infty(\Gamma)\).

\section*{Reference}

M. Gazor, \emph{Condensation rank of injective Banach spaces},
arXiv:1104.4896, 2011.

\end{document}
